% Colors
\definecolor{promptblue}{RGB}{25,55,109}  % Keywords
\definecolor{promptgray}{RGB}{5,51,111}   % Comments

\lstdefinestyle{promptstyle}{
  basicstyle=\ttfamily\small,        % Small monospaced font
  keywordstyle=\color{promptblue}\bfseries, % Bold blue keywords
  commentstyle=\color{promptgray},        % Gray comments
  breaklines=true,                        % Wrap long lines
  breakatwhitespace=true,                 % Break at spaces
  columns=fullflexible,                   % Better spacing
  keepspaces=true,                        % Preserve indentation
  frame=single,                           % Border around code
  framerule=0.4pt,                        % Border thickness
  rulecolor=\color{gray!60},              % Border color
  backgroundcolor=\color{blue!5},         % Light blue background
  numbers=left,                           % Line numbers on left
  numberstyle=\tiny\color{brown},         % Small brown line numbers
  numbersep=8pt,                          % Space from numbers to code
  xleftmargin=30pt,                       % Left margin
  xrightmargin=30pt,                      % Right margin
  aboveskip=15pt,                         % Space above
  belowskip=4pt,                          % Space below
}

% =============================================================
% Appendix: Canonical Prompt Operators
% =============================================================

\section{Prompts}
\label{app:prompts}

This appendix presents canonical prompt templates with formal operators from Algorithm~\ref{alg:rcp_loop} used in the goal–strategy–function control loop.
Each template corresponds to a typed state transition in the RCP and persists structured outputs into the underlying database schema.
Runtime placeholders appear as \texttt{\{variable\_name\}}.

% -------------------------------------------------------------
\subsection{Goal Definition (Stage 1)}

Stage 1 implements the mapping
\( f_G : Q \rightarrow G \),
where \( G \) satisfies schema constraint \( \mathcal{S}_G \) and is persisted in the \texttt{Goal} table.

\noindent\begin{minipage}{\columnwidth}
\begin{lstlisting}[style=promptstyle]
SYSTEM:|
    Expert Goal-definition assistant for user queries.
    - Extract a structured goal from the user query.
    - Define the validation metadata for the goal structure.
    
    Required JSON:
    {"goal_description":"user's technical intent",
     "expected_content_types":["e.g., product_specs, lookup_values, technical_measurements etc"],
     "key_terms": ["domain-specific technical terms"],
     "success_indicators": ["criteria to validate a satisfactory answer"]}
    
USER:|
    USER QUERY: {query}
\end{lstlisting}
\end{minipage}

% -------------------------------------------------------------
\subsection{Strategy Selection (Stage 2)}

Stage 2 implements
\( f_S : (Q, G, \mathcal{L}_S, F) \rightarrow S \),
where \( \mathcal{L}_S \) is the Strategy Library and \( F \) is the set of previously attempted strategies.
Constraint: \( S \in \mathcal{L}_S \setminus F \).
The selected strategy is persisted in the \texttt{Strategy} table.

\noindent\begin{minipage}{\columnwidth}
\begin{lstlisting}[style=promptstyle]
SYSTEM:|
    Strategy planner for a technical documentation system.
    
    - Select EXACTLY ONE strategy from AVAILABLE STRATEGIES
    - Use the EXACT case-sensitive name
    - Do NOT select any strategy listed under FORBIDDEN.
    
    Return valid JSON only:
    {"StrategyName":"[EXACT name]",
     "StrategyTarget":"deep search",
     "Rationale":"BRIEF justification"}

USER:|
    USER QUERY: {query}
    CURRENT GOAL: {goal_desc}
    FORBIDDEN: {tried_readable}
    AVAILABLE STRATEGIES: {lib_block}
\end{lstlisting}
\end{minipage}

% -------------------------------------------------------------
\subsection{Evidence-Constrained Synthesis (Stages 3--5)}

Stages 3--5 implement
\( f_A : (E, Q) \rightarrow A \),
where \( E \) is the assembled evidence set.
Constraint: \( A \subseteq E \) (no unsupported claims).
Outputs are logged in the \texttt{FunctionOutput} and execution-trace tables.

\noindent\begin{minipage}{\columnwidth}
\begin{lstlisting}[style=promptstyle]
SYSTEM:|
    EXPERT technical data analyst. 
    Analyze the compiled data and provide clear, precise answers to user queries, ONLY with the provided DATA CONTEXT. If information is missing, state so explicitly.
    
    - Cross-reference codes with spec tables
    - Retrieve via linking fields
    - Apply semantic field mapping
    - Return exact specifications
    
    Field mappings:
    {- contact count -> Number of contacts # size
    - shell size -> Shell size
    - cable diameter -> Max Cable entry
    - torque -> fields with "torque" or "Nm"}

USER:|
    DATA CONTEXT: {combined_context}
    USER QUERY: {query}
\end{lstlisting}
\end{minipage}

% -------------------------------------------------------------
\subsection{Goal Validation (Stage 6)}

Stage 6 implements the acceptance predicate
\( f_V : (G, A) \rightarrow c \in [0,1] \),
where \( c \) is a confidence score used by the RCP termination logic.
If \( c \geq \tau \), the goal is marked successful.

\noindent\begin{minipage}{\columnwidth}
\begin{lstlisting}[style=promptstyle]
SYSTEM:|
    You are a goal-level evaluator for technical queries.
    Return ONLY valid JSON in the form:
    {"confidence": 0.0}
    
    Confidence Scoring:
    - 0.8-1.0 = complete, detailed answer
    - 0.6-0.7 = good coverage of main aspects
    - 0.3-0.5 = partial / insufficient
    - 0.0-0.2 = no meaningful answer

USER:|
    USER QUERY: {query}
    GOAL DEFINITION: {goal_definition}
    FINAL OUTPUT: {full_evidence}
\end{lstlisting}
\end{minipage}

\medskip
\noindent
All prompt inputs and structured outputs are persisted in the RCP execution-trace schema, enabling deterministic replay, failure analysis, and post-hoc validation.
